\section{Discussion}

The empirical results across TwiBot-20, TwiBot-22, and Cresci-2017 strongly validate the core hypothesis of the XHBot framework: standard homophilic message passing is fundamentally insufficient for modern bot detection due to adversarial relation camouflage.

\subsection{Scalability and Cross-Domain Generalization}
A major challenge in social bot detection is scale and graph heterogeneity. The performance of XHBot on the massive TwiBot-22 benchmark (which contains approximately 1 million nodes and 170 million edges) demonstrates that the Tri-Channel Heterophily-Aware Aggregation (THCA) mechanism not only resists camouflage but also scales efficiently to highly complex, multi-relational schemas without catastrophic performance degradation. While foundational models like BotRGCN collapse to 0.68 F1 on TwiBot-22 due to noise dilution, XHBot maintains a highly competitive 0.9125 F1. Furthermore, the near-perfect 0.9890 F1 score on Cresci-2017 confirms that XHBot easily captures the simpler structural signatures of older spambot generations, proving its backward compatibility.

\subsection{The Value of Forensic Transparency}
Beyond classification accuracy, XHBot is uniquely positioned for real-world deployment due to its Hierarchical Forensic Explanation (HFE) module. Our quantitative explainability evaluations confirm high \textbf{Fidelity+} (a significant drop in model confidence when the explanation subgraph is masked) and high \textbf{Fidelity-} (maintained accuracy when only the explanation subgraph is preserved), alongside high \textbf{Sparsity} (generating concise, human-readable subgraphs). 

By combining instance-level GNNExplainer masks with Louvain community detection, platform moderators are provided with actionable, localized proof of botnet coordination—shifting the paradigm from black-box probability scoring to transparent, evidence-based moderation. This is critical for complying with emerging AI transparency regulations and maintaining user trust.
